\documentclass[11pt,a4paper]{article}
\usepackage[T1]{fontenc}
\usepackage[margin=24mm]{geometry}
\usepackage{amsmath,amssymb,amsthm,mathtools}
\usepackage[hidelinks]{hyperref}
\allowdisplaybreaks[2]
\setlength{\emergencystretch}{3em}
\newtheorem{theorem}{Theorem}[section]
\newtheorem{lemma}[theorem]{Lemma}
\newtheorem{corollary}[theorem]{Corollary}
\DeclareMathOperator{\arsinh}{arsinh}
\DeclareMathOperator{\im}{im}
\DeclareMathOperator{\rank}{rank}
\DeclareMathOperator{\Gr}{Gram}
\newcommand{\R}{\mathbb R}
\newcommand{\C}{\mathbb C}
\newcommand{\norm}[1]{\left\lVert#1\right\rVert}
\renewcommand{\theequation}{QOT\arabic{equation}}
\title{Native Gamma-order transport of the common long loss\\and its low primitive source}
\author{Split-Zero research continuation}
\date{16 September 2026}
\begin{document}
\maketitle
\begin{abstract}
For the original fixed simple-quartet family, a full quadratic-form comparison controls the change from Gamma order one to the original degree-dependent order. It acts on every original coefficient direction before either relation or observation minimization. The difference of the two native kernel allocations, and separately of the two residual allocations, is $O(q_k)$, with limiting absolute allowance $8\log(2+\sqrt3)$ after division by $q_k$. The common connecting/invariant long loss has source-order difference $O(q_k)$ with allowance $32\log(2+\sqrt3)$. At the two defining primitive cutoffs, the paired source-to-connecting-image order correction is $o(q_k)$. Its scalar shift is evaluated exactly, and the original compatibility difference is retained and compared in the same coordinates. The individual $kq_k$-scale allocation coefficients and the absolute common long loss are not evaluated here. No offcritical packet is excluded by these comparisons alone.
\end{abstract}
% BEGIN EOR RECEIVER EDITION
\paragraph{Refinement retained in this edition.}
The subsequent EOR calculation, included in full below, replaces the high-degree
symmetric enclosure by the actual equilibrium-support enclosure. For
$\Delta S_X=S_X(k;k)-S_X(k;1)$ it gives the limiting interval
$[-8\log(\pi v_*/4),-8\log(\pi u_*/4)]$ at scale $q_k$.
The common source-order allowance becomes $16\log(v_*/u_*)$.
The original QOT estimates remain valid; these are subsequent strict refinements,
not a reattribution of the older results. The original source bytes are preserved.
Two truncated norm control sequences in its source have been restored here,
with an exact reversible record.
% END EOR RECEIVER EDITION

\tableofcontents

\section{Original data and the result's scope}
Fix the stipulated hypothetical actual simple quartet, its full unit, one original five-orbit period and its logarithm branch. Put
\begin{equation}
\begin{gathered}
 k=4l+1\ge9,\qquad q_k=(k+1)^2,\quad q'_k=(k-7)^2,\quad
 \Delta_k=16k-48,\quad c'_k=k/2-4,\\
 \beta_{n,a,b}=n/2+(2a-n)\delta+i(2b-n)\gamma,
 \quad0<\delta<1/2,\quad\gamma>2,\\
 \chi_n(S)=\prod_{a,b=0}^n(S-\beta_{n,a,b}),\qquad \psi_k=\chi_{k-8},\\
 E_A(z)=\sum_{a,b=0}^8a_{ab}e^{\beta_{8,a,b}z},\quad
 v=\operatorname{ord}_0E_A,\quad \mu_v=E_A^{(v)}(0)\ne0,\quad
 \alpha_A=v!/\mu_v .
\end{gathered}\label{eq:data}
\end{equation}
The $a_{ab}$ are the actual conductor coefficients. Its nonreduced factors and exact $v$ remain. The fixed bound $v\le80$ is sufficient; no exact matrix replaces $v$ by that bound. The polynomial $\psi_k$ is a lower-grid polynomial, not a derivative.

At degree $k$ the two original source orders are $s=1,k$. The complete measures, masses, and monic norms are
\begin{equation}
\begin{gathered}
 dm_s(y)=\frac{M_s2^{s/2}}{4\pi\Gamma(s/2)}
     |\Gamma(s/4+iy/2)|^2dy,\qquad M_s=(2\pi)^{s/2},\\
 h_{n,s}=M_sn!(s/2)_n,\quad p_{0,s}=1,\quad p_{1,s}=y,\\
 p_{n+1,s}=yp_{n,s}-n(n-1+s/2)p_{n-1,s}.
\end{gathered}\label{eq:measure}
\end{equation}
All comparisons below identify the same original polynomial coefficients or quotient values, not their different orthonormal coordinate vectors. Source orders are never altered inside an endpoint telescoping sum.

The quantitative root/Gamma input used below is the supplied full-form GPA proof. Its original finite domain is
\begin{equation}
 k\ge29,\qquad R_k=k\sqrt{\delta^2+\gamma^2}\le2^{-33}q_k .
 \label{eq:domain}
\end{equation}
For paired degree $j=k+4$ statements impose the corresponding domain at $j$ as well. Every fixed actual quartet eventually satisfies both conditions. Arithmetic statements retain the original period domain $|u|\ge R_*$, or the specified original exact period tests. There is no uniform assertion over a varying period.

The available native primal forms are the attained minima
\begin{equation}
\begin{split}
 z^*H_{Y,k,N}^{(s)}z
 &=\min_h\int_{\R}|(Ez+\mathsf U_Nh)(c'_k+iy)|^2
                   |\psi_k(c'_k+iy)|^2dm_s(y),\\
 H_{K,k,N}^{(s)}&=Z_k^*H_{Y,k,N}^{(s)}Z_k,\\
 H_{R,k,N}^{(s)}&=[P_Y(H_{Y,k,N}^{(s)})^{-1}P_Y^*]^{-1}.
\end{split}\label{eq:primal}
\end{equation}
Here $E$, $\mathsf U_N$, $Z_k$, and $P_Y$ are the complete original GPA/FGC/SAR maps. Each relation column is $\alpha_A\mathcal T_A(\chi_k S^d)/\psi_k$, where $\mathcal T_Ap(S')=\sum a_{ab}p(S'+\beta_{8,a,b})$. Their domains, all coefficients and both minima in \eqref{eq:primal} remain unchanged.
The ranks and endpoint signs are
\begin{equation}
\begin{gathered}
 \mathfrak m_k=\Delta_k-v,\quad t_0=\dim(K\cap\C[S]_{<v}),\quad
 t_K=8k-16-t_0,\quad t_R=8k-32-v+t_0,\\
 N=(q_k-1,q_k,2q_k-1,2q_k),\quad \sigma=(+1,+1,-1,-1),\\
 S_{X,k,s}=\sum_N\sigma_N\log\det H_{X,k,N}^{(s)},\qquad X=K,R.
\end{gathered}\label{eq:alloc}
\end{equation}
The actual dual coefficient metric is
\begin{equation}
 D_{k,N-v}^{(s)}=A_{\ell,k}^*\overline{(H_{Y,k,N}^{(s)})^{-1}}A_{\ell,k}.
 \label{eq:dual}
\end{equation}
The same invertible, source-order-independent $A_{\ell,k}$ appears in both metrics. The conjugation is required by complex-linear duality.

\section{A complete power-weight comparison}
For $\alpha>0$, $\vartheta>0$, let
\[
 \norm{p}_\alpha^2=\int_{\R}|p(y)|^2|y|^\alpha e^{-\vartheta|y|}dy.
\]
\begin{lemma}[Multiplication with its next-degree row retained]\label{lem:mult}
For $n\ge0$ define
\begin{equation}
\begin{split}
 c_n&=\frac{(\sqrt{n+\alpha+1/2}-\sqrt{n+1/2})^2}{\vartheta},\\
 C_n&=\frac{(\sqrt{n+\alpha+1/2}+\sqrt{n+1/2})^2}{\vartheta}.
\end{split}\label{eq:cn}
\end{equation}
Every complex polynomial of degree at most $n$ satisfies
\begin{equation}
 c_n\norm{p}_\alpha\le\norm{yp}_\alpha\le C_{n+1}\norm{p}_\alpha.
 \label{eq:mult}
\end{equation}
\end{lemma}
\begin{proof}
On $(0,\infty)$, the normalized Laguerre polynomials for $y^\alpha e^{-\vartheta y}dy$ give multiplication matrix
\[
 (J_n)_{ii}=(2i+\alpha+1)/\vartheta,\qquad
 (J_n)_{i,i+1}=-\sqrt{(i+1)(i+\alpha+1)}/\vartheta
 \quad(0\le i<n).
\]
Its entries follow from Rodrigues' coefficient formula and integration by parts, as in the supplied GPA/LVM proof. For completeness, $L_n^{(\alpha)}(x)=\sum_{r=0}^n(-1)^r\binom{n+\alpha}{n-r}x^r/r!$ has squared norm $\Gamma(n+\alpha+1)/n!$ and obeys
\[
 xL_n^{(\alpha)}=(2n+\alpha+1)L_n^{(\alpha)}
 -(n+1)L_{n+1}^{(\alpha)}-(n+\alpha)L_{n-1}^{(\alpha)}.
\]
Coefficient comparison proves the recurrence; the stated norms give $J_n$, including its negative off-diagonal entries.

The function $f(x)=\sqrt{x(x+\alpha)}$ is concave for $x>0$, since $f''(x)=-\alpha^2/[4(x(x+\alpha))^{3/2}]$. Continuity includes $x=0$. Hence $f(i)+f(i+1)\le2f(i+1/2)$. Bounding the cross terms of the tridiagonal quadratic form by the adjacent squared moduli gives
\[
 c_n I\preceq J_n\preceq C_n I.
\]
Indeed the lower row expression $2i+\alpha+1-2f(i+1/2)$ decreases with $i$, and the upper expression increases. Omitting a missing boundary off-diagonal only improves these bounds. Cauchy--Schwarz now gives $\norm{yp}\norm p\ge\langle yp,p\rangle\ge c_n\norm p^2$ on this half-line.

For the upper norm bound the multiplication map is rectangular: its last row is $-\sqrt{(n+1)(n+\alpha+1)}\,e_n^{\mathsf T}/\vartheta$. It is a restriction of $J_{n+1}$, not $J_n$. Thus its norm is at most $C_{n+1}$. Apply both inequalities also to $p(-y)$ on $y>0$ and add their squares. The sign of multiplication on the negative half-line has no effect on its norm. This proves \eqref{eq:mult} for the entire original coefficient direction, without a parity restriction.
\end{proof}

For integers $a,\ell\ge0$ put
\begin{equation}
\begin{split}
 d^-_{\alpha,a,\ell}
 &=4\sum_{r=0}^{\ell-1}\arsinh\sqrt{(a+r+1/2)/\alpha},\\
 d^+_{\alpha,a,\ell}
 &=4\sum_{r=1}^{\ell}\arsinh\sqrt{(a+r+1/2)/\alpha}.
\end{split}\label{eq:ds}
\end{equation}
Empty sums are zero.
\begin{theorem}[Centred weight shift]\label{thm:weights}
On every polynomial of degree at most $a$,
\begin{equation}
 (\alpha/\vartheta)^{2\ell}e^{-d^-_{\alpha,a,\ell}}\norm p_\alpha^2
 \le\norm p_{\alpha+2\ell}^2
 \le(\alpha/\vartheta)^{2\ell}e^{d^+_{\alpha,a,\ell}}\norm p_\alpha^2.
 \label{eq:weightshift}
\end{equation}
\end{theorem}
\begin{proof}
Iterate \eqref{eq:mult} in the fixed $\alpha$ norm on $p,yp,\ldots,y^{\ell-1}p$, with degree bounds $a,a+1,\ldots,a+\ell-1$. The resulting lower product is $\prod_{r=0}^{\ell-1}c_{a+r}^2$ and the upper product is $\prod_{r=1}^{\ell}C_{a+r}^2$. Also $\norm{y^\ell p}_\alpha^2=\norm p_{\alpha+2\ell}^2$ exactly. Finally
\[
 c_n=(\alpha/\vartheta)e^{-2\arsinh\sqrt{(n+1/2)/\alpha}},\quad
 C_n=(\alpha/\vartheta)e^{2\arsinh\sqrt{(n+1/2)/\alpha}}.
\]
Substitution proves both sides. The displayed central scalar is not an exact equality of forms; for $p=1$, $\ell=1$, the actual moment ratio is $(\alpha+1)(\alpha+2)/\vartheta^2$, not $\alpha^2/\vartheta^2$.
\end{proof}

\section{Finite constants in the original Gamma comparison}
We record the full inherited GPA constants so the comparison does not use an unspecified conditioning constant. Abbreviate $q=q_k$, $q'=q'_k$, $R=R_k$, $\epsilon=2^{-32}$, $\vartheta=\pi/2$, and $h=1/q$. Put
\begin{equation}
\begin{gathered}
 \ell_s=(s-1)/4,\quad Q_s=q'+\ell_s,\quad
 \beta_s=\frac{(2\pi)^{s/2}}{\sqrt2\,\Gamma(s/2)},\quad \beta_1=1,\\
 C_L=\frac\pi{\Gamma(3/4)^2},\quad
 K_q=\frac{\sqrt{2\pi}}{C_L}\frac{(2q+1)^2}{q}e^{-2q},\\
 \mathcal E_k=\frac{qR^2}{\epsilon^2q^2-R^2}+\log(1+K_q),\\
 r_{\ell_s}=\frac{\sum_{r=0}^{\ell_s-1}(2r+1/2)^2}{\epsilon^2q^2}
                 +\log(1+K_q)\quad(\ell_s>0),\qquad r_0=0,\\
 L_h=C_L(\sin h)^{3/2},\quad
 U_h=\frac{\Gamma(1/4)^2}{2\pi}(\sin h)^{-1/2},\\
 K_{a,Q}=2Q+1+2\{a+\sqrt{a(a+2Q)}\}.
\end{gathered}\label{eq:GPconst}
\end{equation}
The logarithmic endpoints are
\begin{equation}
\begin{split}
 b^-_{k,a,s}&=-\mathcal E_k+\log L_h-K_{a,Q_s}\log(1+h/\vartheta),\\
 b^+_{k,a,s}&=\mathcal E_k+r_{\ell_s}+\log U_h
                    +K_{a,Q_s}\log\frac1{1-h/\vartheta}.
\end{split}\label{eq:bconst}
\end{equation}
The supplied full proof gives, for every $\deg p\le a\le2q$,
\begin{equation}
\begin{split}
 \beta_s e^{b^-_{k,a,s}}
\norm{p(c'_k+i\,\cdot)}_{2Q_s}^2
 &\le \int_{\R}|\psi_k(c'_k+iy)p(c'_k+iy)|^2dm_s(y)\\
 &\le\beta_s e^{b^+_{k,a,s}}
\norm{p(c'_k+i\,\cdot)}_{2Q_s}^2.
\end{split}\label{eq:GPA}
\end{equation}
This uses precisely the degree bound in the proof of GPA12--15, not only its originally listed endpoints. In all applications here $a+\ell_s\le2q$ also holds on the stated eventual domains. The proof pairs the original opposite roots outside $[-\epsilon q,\epsilon q]$, bounds the entire inner polynomial integral by $K_q$ times the reference outside integral, and retains the full source recurrence
\[
 dm_s=\beta_s\prod_{r=0}^{\ell_s-1}(y^2+(2r+1/2)^2)d\sigma,
 \qquad d\sigma=dm_1.
\]
It then uses the whole-line Gamma sandwich $L_he^{-(\vartheta+h)|y|}\le d\sigma/dy\le U_he^{-(\vartheta-h)|y|}$ and the exact polynomial dilation norm bound. Thus none of the original roots, the inner region, or the source mass is removed by \eqref{eq:GPA}. At fixed actual data, $b^\pm_{k,a,s}=O_{\rm actual}(\log q)$ uniformly for these degrees and both original orders.

Define
\begin{equation}
\begin{gathered}
 \alpha_k=2q'_k,\quad \ell=(k-1)/4,\qquad
 \kappa_k=\beta_k(4q'_k/\pi)^{(k-1)/2},\\
 l_{k,a}=b^-_{k,a,k}-b^+_{k,a,1}-d^-_{\alpha_k,a,\ell},\qquad
 u_{k,a}=b^+_{k,a,k}-b^-_{k,a,1}+d^+_{\alpha_k,a,\ell},\qquad
 w_{k,a}=u_{k,a}-l_{k,a}\ge0.
\end{gathered}\label{eq:ordconst}
\end{equation}
\begin{theorem}[Full original source-order transport]\label{thm:full}
On the degree-$a$ original coefficient space, and on every fixed restriction or attained quotient of it,
\begin{equation}
 \kappa_k e^{l_{k,a}}H^{(1)}\preceq H^{(k)}
                 \preceq\kappa_k e^{u_{k,a}}H^{(1)}.
 \label{eq:primalord}
\end{equation}
In particular it holds for all three forms in \eqref{eq:primal} at $a=N-v-q'_k$. On their original complex-linear dual coefficient metric it becomes
\begin{equation}
 \kappa_k^{-1}e^{-u_{k,a}}D^{(1)}\preceq D^{(k)}
                 \preceq\kappa_k^{-1}e^{-l_{k,a}}D^{(1)}.
 \label{eq:dualord}
\end{equation}
The same dual inequality passes to every specified fixed flag subquotient, including the fourfold target space.
\end{theorem}
\begin{proof}
Apply \eqref{eq:weightshift} with $\alpha=2q'$ and $\ell=(k-1)/4$ to the complete $p(c'+iy)$ direction. Combine its lower bound with the lower order-$k$ inequality and upper order-one inequality in \eqref{eq:GPA}; combine the three opposite inequalities for the upper side. This gives exactly \eqref{eq:ordconst}--\eqref{eq:primalord}. On a restriction substitute the same coefficient columns. On a quotient every candidate in the same affine fibre obeys both inequalities, so its minimum does as well. This preserves the entire relation and kernel Schur corrections. Inversion reverses positive-form inequalities, complex conjugation preserves them, and the fixed congruence \eqref{eq:dual} proves \eqref{eq:dualord}. Apply the same restriction/minimum argument to its subsequent fixed flags.
\end{proof}

\section{The two-half-line transport is literal}
To specify the coefficient comparison completely, for $\alpha=2Q_s$ and $0\le j\le r\le a$ let
\begin{equation}
 \mathsf A_{j,r}=(-1)^j\binom rj\Gamma(r+\alpha+1)
 \sqrt{\frac{j!}{\Gamma(j+\alpha+1)}}
 \vartheta^{-r-(\alpha+1)/2},\qquad
 \mathsf P_{r,b}=\binom br(c'_k)^{b-r}i^r .
 \label{eq:rowsAP}
\end{equation}
Other triangular entries are zero. These are respectively the monomial-to-orthonormal-Laguerre and original affine coefficient maps. The exact reflection is
\begin{equation}
 \mathsf R_{j,r}=(-1)^j2^{r-j}\binom{\alpha+r}{r-j}
  \sqrt{\frac{r!\Gamma(j+\alpha+1)}{j!\Gamma(r+\alpha+1)}},
 \qquad \mathsf R^2=I,
 \qquad \mathsf W=\begin{bmatrix}\mathsf A\mathsf P\\
                    \mathsf R\mathsf A\mathsf P\end{bmatrix}.
 \label{eq:W}
\end{equation}
The Gram $\mathsf W^*\mathsf W$ is the complete reference integral over both half-lines. Reflection is complex-linear: it does not conjugate a period coefficient. For any actual relation or invariant column $x$, the two row blocks are $(\mathsf A\mathsf P)x$ and $(\mathsf R\mathsf A\mathsf P)x$. The source change in Theorem~\ref{thm:full} is on these full integrals before minimization; it is not a separate fit on each half-line. In particular the same $\mathsf U_N$, $Z_k$, and $P_Y$ enter their two quadratic forms. All original factorials and the source factor $\beta_s$ have been displayed, rather than assigned to an unknown inverse.

\section{An evaluated allowance at scale $q_k$}
For $X=K,R$ let $a_N=N-v-q'_k$. Theorem~\ref{thm:full} and the four original signs give
\begin{equation}
\begin{split}
 t_X(l_{k,a_{q-1}}+l_{k,a_q}-u_{k,a_{2q-1}}-u_{k,a_{2q}})
 &\le S_{X,k,k}-S_{X,k,1}\\
 &\le t_X(u_{k,a_{q-1}}+u_{k,a_q}-l_{k,a_{2q-1}}-l_{k,a_{2q}}).
\end{split}\label{eq:allocationinterval}
\end{equation}
The four copies of $t_X\log\kappa_k$ cancel only in this signed sum.

In fact the scalar is an explicit positive rational number at every original degree:
\begin{equation}
 \boxed{\qquad
 \kappa_{4l+1}=2^{10l}\frac{(2l)!}{(4l)!}(k-7)^{4l}.
 \qquad}\label{eq:kapparational}
\end{equation}
This follows from $\Gamma(2l+1/2)=\sqrt\pi(4l)!/[4^{2l}(2l)!]$.
The powers of $\pi$ in the original mass and the central power weight
cancel only in this exact scalar identity.


At a low endpoint, $a=\Delta_k-v-1$ or $\Delta_k-v$ is $O(k)$. Since $\arsinh x\le x$ for $x\ge0$, \eqref{eq:ds} gives $d^\pm=O(\sqrt k)$, hence $l,u=O_{\rm actual}(\sqrt k+\log q_k)$. At either high endpoint, $a=q_k+\Delta_k-v+O(1)$ and $\alpha_k=2q'_k$. Uniformly over the $\ell=O(k)$ summands,
\[
 \frac{a+r+1/2}{\alpha_k}\longrightarrow\frac12.
\]
Therefore
\begin{equation}
 \frac{l_{k,a}}k\longrightarrow-a_*,\qquad
 \frac{u_{k,a}}k\longrightarrow a_*,\qquad
 a_*=\arsinh(1/\sqrt2)=\tfrac12\log(2+\sqrt3).
 \label{eq:limits}
\end{equation}
Since $t_X/k\to8$ and $q_k/k^2\to1$, \eqref{eq:allocationinterval} proves
\begin{equation}
 \boxed{\quad
 \limsup_{k\to\infty}\frac{|S_{X,k,k}-S_{X,k,1}|}{q_k}
 \le 8\log(2+\sqrt3),\qquad X=K,R .\quad}
 \label{eq:allocationq}
\end{equation}
% BEGIN EOR RECEIVER ALLOCATION
\paragraph{Sharper receiving interval (EOR24--26).}
In the same original source and quotient coordinates, use the finite endpoints
\eqref{eor:eq:allocationfinite} wherever their guards hold, intersected as in
\eqref{eor:eq:intersection}. The resulting limiting interval is
\[\begin{gathered}
 -C_+\le\liminf\frac{S_{X,k,k}-S_{X,k,1}}{q_k}
 \le\limsup\frac{S_{X,k,k}-S_{X,k,1}}{q_k}\le C_-,\\
 C_+=8\log(\pi v_*/4),\qquad C_-=8\log(4/(\pi u_*)).
\end{gathered}\]
Here $C_+<8.881817865597$ and $C_-<5.773751910974$; neither is an individual
allocation coefficient. The full four endpoint scalars still cancel with their
original signs.
% END EOR RECEIVER ALLOCATION

This is an allowance for the difference of the two source-order allocations, not the asymptotic coefficient of either allocation. In particular their $kq_k$-normalized limit points coincide.

There is a genuine evaluated low-endpoint order-change term. Stirling's elementary factorial bounds give
\begin{equation}
 \log\kappa_k=\tfrac{k}{2}\log k+\tfrac{k}{2}(1+4\log2)+O(\log k).
 \label{eq:kappaasym}
\end{equation}
It can also be obtained from the integer/half-integer Gamma formula and integral bounds for $\sum\log n$. At either low endpoint, Theorem~\ref{thm:full} gives an error $O_{\rm actual}(k^{3/2}+k\log q_k)=o(q_k)$ after taking a rank-$t_X$ determinant. Thus
\begin{equation}
\boxed{\begin{gathered}
 \log\frac{\det H_{X,k,q_k-1+e}^{(k)}}{
                 \det H_{X,k,q_k-1+e}^{(1)}}
 =4q_k\log k+(4+16\log2)q_k+o(q_k),\\
 e=0,1,\qquad X=K,R.
\end{gathered}}
\label{eq:lowprimal}
\end{equation}
No next-order term for the complete four-return is assigned: its high determinants still require their actual allocation calculation.

\section{The defining primitive source and connecting image}
Put $j=k+4$, $Q=q_j$, $D_e=Q-v-1+e$, and $d_e=8k+24+e$. The original PDS relations are
\[
 r_a=\mathcal T_A(\chi_k S^a),\qquad 0\le a<d_e,
 \qquad \deg r_a=q_k+a-v.
\]
\begin{lemma}[The original lower factor divides every primitive column]\label{lem:divide}
Every $r_a$ is divisible by the full $\psi_k$. Its quotient degree is $\Delta_k-v+a$.
\end{lemma}
\begin{proof}
For a root $\beta_{k-8,u,w}$ of $\psi_k$, every summand in $r_a(\beta_{k-8,u,w})$ evaluates $\chi_k$ at
$\beta_{k-8,u,w}+\beta_{8,b,c}=\beta_{k,u+b,w+c}$, one of its original upper roots. Each summand is zero, without a cancellation or support assumption. The lower roots are distinct, so their complete product divides $r_a$. The degree follows from the first nonzero conductor moment $\mu_v$.
\end{proof}
In particular the whole relation Gram $H_e$ is a restriction of the original weighted form in Theorem~\ref{thm:full} with degree bound
\begin{equation}
 a_s=\Delta_k-v+d_e-1=24k-25-v+e .
 \label{eq:asrc}
\end{equation}
This bound is within $2q_k$ on the stated domain. No leading-monomial approximation or unproved inverse-conductor estimate enters.

We next identify the same primitive coordinates in both source orders. Let $R_e^{\rm pol}$ be the full coefficient matrix of the $r_a$, let $\mathsf B_s$ be the monomial-to-original-Gamma coefficient matrix, and let $P_e^{\rm Tay}$ be CIT's fixed Taylor complement. Then
\begin{equation}
\begin{gathered}
 U_s=\operatorname{diag}(1/r!)\mathsf B_s^{\mathsf T},\quad
 P_e^{(s)}=U_s^{-1}P_e^{\rm Tay},\quad C_e^{(s)}=\mathsf B_sR_e^{\rm pol},\\
 \Theta_e=(C_e^{(s)})^{\mathsf T}P_e^{(s)}
 =(R_e^{\rm pol})^{\mathsf T}\operatorname{diag}(r!)P_e^{\rm Tay}.
\end{gathered}\label{eq:Theta}
\end{equation}
The actual $\Theta_e$ is independent of $s$; every mass, affine phase, and factorial cancels here against its exact inverse, not against an assumed norm. At $e=0$ it is invertible; at $e=1$ it has image the same rank-$r_e$ primitive constraint space, where $r_e=d_e-4e=8k+24-3e$. PDS/CIT's equality of affine fibres gives
\[
 H_{{\rm src},e}^{(s)}=\Theta_e^*\overline{(H_e^{(s)})^{-1}}\Theta_e.
\]
Put $l_s=l_{k,a_s}$, $u_s=u_{k,a_s}$. Theorem~\ref{thm:full} therefore proves
\begin{equation}
 -r_eu_s\le
 \log\frac{\det H_{{\rm src},e}^{(k)}}{\det H_{{\rm src},e}^{(1)}}
       +r_e\log\kappa_k
 \le-r_el_s.
 \label{eq:sourceorder}
\end{equation}
This restriction retains all four original primitive constraints at $e=1$.

The corresponding CIT image columns are also the same coefficient columns in both source pairs. In ordinary Taylor coordinates their equation is
\[
 \operatorname{stack}_i\operatorname{Mult}(B_i)P_e^{\rm Tay}
 =\operatorname{diag}(\widetilde T_j,4)F_e
   +\operatorname{diag}(\widetilde V_j,4)E_e^{\rm low}.
\]
All Taylor matrices are those of the same specified germs. The original analytic quotient is injective at this target cutoff modulo the full lower evaluations, so its coefficient $F_e$ is unique. This proves equality of the two $F_e$ and of $B_{\rm cof}$ in their common value coordinates. It is the exact PSR23--24/CIT coefficient argument, not an inference from equal leading determinant orders.

The connecting Gram is the complete fixed quotient
\begin{equation}
\begin{split}
 H_{{\rm conn},e}(N)={}&F_e^*\mathcal D_NF_e
 -F_e^*\mathcal D_NB_{\rm cof}
 (B_{\rm cof}^*\mathcal D_NB_{\rm cof})^{-1}
 B_{\rm cof}^*\mathcal D_NF_e,\\
 &\mathcal D_N=\operatorname{diag}(D_{j,N-v}^{(s_1)},4).
\end{split}\label{eq:conn}
\end{equation}
At its defining low cutoff $N=Q-1+e$, let
\begin{equation}
 a_t=\Delta_j-v-1+e=16k+15-v+e,
 \qquad l_t=l_{j,a_t},\quad u_t=u_{j,a_t}.
 \label{eq:atgt}
\end{equation}
The dual comparison and the identical complete boundary minimum give
\begin{equation}
 -r_eu_t\le
 \log\frac{\det H_{{\rm conn},e}^{(k,j)}}{
                 \det H_{{\rm conn},e}^{(1,1)}}
       +r_e\log\kappa_j
 \le-r_el_t.
 \label{eq:targetorder}
\end{equation}
All four constants $l_s,u_s,l_t,u_t$ are $O_{\rm actual}(\sqrt k+\log q_k)$.

\subsection{The exact scalar cancellation and its compatibility return}
The Gamma recurrence gives the elementary identity
\begin{equation}
\boxed{\begin{gathered}
 \frac{\kappa_{k+4}}{\kappa_k}
 =\frac{256(k-3)^{k+3}}{k(k+2)(k-7)^{k-1}},\\
 \log\frac{\kappa_{k+4}}{\kappa_k}
 =2\log k+8\log2+4+O(k^{-1}).
\end{gathered}}
\label{eq:kappapair}
\end{equation}
Indeed $\beta_{k+4}/\beta_k=16\pi^2/[k(k+2)]$ and the two remaining powers supply $16/\pi^2$ and the displayed integer powers. Every original source mass is included in this cancellation.

Define the difference of the two original source-to-connecting-image logarithmic Jacobians by
\[
 \mathcal J_e^{\rm ord}=
 \log\frac{\det H_{{\rm conn},e}^{(k,j)}}{\det H_{{\rm src},e}^{(k)}}
 -\log\frac{\det H_{{\rm conn},e}^{(1,1)}}{\det H_{{\rm src},e}^{(1)}}.
\]
Subtracting \eqref{eq:sourceorder} from \eqref{eq:targetorder}, in the same original primitive values, proves
\begin{equation}
\boxed{\begin{gathered}
 r_e(l_s-u_t)\le
 \mathcal J_e^{\rm ord}+r_e\log(\kappa_j/\kappa_k)
 \le r_e(u_s-l_t),\\
 \mathcal J_e^{\rm ord}=o(q_k).
\end{gathered}}
\label{eq:paired}
\end{equation}
The error before division by $q_k$ is $O_{\rm actual}(k^{3/2}+k\log q_k)$. This bound can exceed the displayed $O(k\log k)$ scalar centre; the centre is not asserted to be the next asymptotic term. Separately, the two low order changes in \eqref{eq:sourceorder} and \eqref{eq:targetorder} each equal
\begin{equation}
 -4q_k\log k-(4+16\log2)q_k+o(q_k).
 \label{eq:lowdual}
\end{equation}
Thus the common large term cancels only in their paired Jacobian. Neither absolute low Gram is declared close to its order-one value.

Write $\mathcal X_e^{(s_0,s_1)}=\chi_{X_k,D_e}-\chi_{Y_e,D_e}$ for CIT's full proper-source compatibility difference. Its exact identity is
\[
 \log\det H_{{\rm conn},e}-\log\det H_{{\rm src},e}
 =-\mathcal X_e+\varepsilon_e,
\]
where the original finite endpoints are
\[
 \varepsilon_e^-=-2r_e\log\max(1,K_-)-r_e\log c_+,
 \qquad
 \varepsilon_e^+=2r_e\log\max(1,K_+)-r_e\log c_-.
\]
The constants $K_\pm,c_\pm$ are exactly the original multiplier and compatibility constants at that cutoff and source pair. They are not redefined by the present comparison; their supplied bounds are $O_{\rm actual}(k\log q_k)$ after the displayed rank factors.
Let $\Delta\mathcal X_e=\mathcal X_e^{(k,j)}-\mathcal X_e^{(1,1)}$. Then
\begin{equation}
\begin{split}
 r_e(l_t-u_s)+\varepsilon_{e,(k,j)}^-
                  -\varepsilon_{e,(1,1)}^+
 &\le\Delta\mathcal X_e-r_e\log(\kappa_j/\kappa_k)\\
 &\le r_e(u_t-l_s)+\varepsilon_{e,(k,j)}^+
                  -\varepsilon_{e,(1,1)}^- .
\end{split}\label{eq:chi}
\end{equation}
In particular $\Delta\mathcal X_e=o(q_k)$. Only the difference between source-order pairs is small. The absolute compatibility loss is not bounded here. At $e=0$ the original $\chi_Y$ is zero by its proved full-source statement; at $e=1$ all four constraints and the complete $\chi_Y$ remain.

\section{Source-order control for the original common long loss}
At every target endpoint use the fixed coefficient flags
\[
 I=\im(B_{\rm cof}J_0),\quad B=\im B_{\rm cof},\quad
 S=\im[B_{\rm cof},F_e],\quad J=\ker\mathbb B,
\]
\[
 K_i=I\cap J,\quad W_i=I+J,\quad W=B+J,\quad
 K_c=S\cap W,\quad T=S+J.
\]
The original losses in the common value frames are
\begin{equation}
 \omega_i=\frac{\det\Gr(W_i/J)}{\det\Gr(I/K_i)},\qquad
 \omega_c=\frac{\det\Gr(T/W)}{\det\Gr(S/K_c)}.
 \label{eq:loss}
\end{equation}
These are ITL's invariant and CIT's connecting losses. Put $r_i=\dim(I/K_i)$, $r_c=\dim(S/K_c)$, and $R_e=r_i+r_c$. The complete kernel of the original multiplication/observation map is retained in $K_i,K_c$. In particular
\[
 R_e\le(8k-32-v+t_0)+(8k+24-3e)\le16k-8.
\]
All the flags are literally independent of source order, by their polynomial maps and the uniqueness argument preceding \eqref{eq:conn}.

A general immediate consequence of \eqref{eq:dualord} is useful. If two matched fixed subquotients have the same rank $r$, their determinant ratio obeys
\begin{equation}
 \left|\log\omega^{(j)}(N)-\log\omega^{(1)}(N)\right|
 \le r\,w_{j,N-v-q'_j}.
 \label{eq:ratcompare}
\end{equation}
Indeed the logarithm of each determinant's order change belongs to $[-r\log\kappa_j-ru,-r\log\kappa_j-rl]$. Subtracting these intervals cancels the scalar before bounding. This proof does not posit a lower singular-value gap for an observation map.

Let $\mathcal R$ be the original native-dual return at $N=q_j-1,q_j,2q_j-1,2q_j$ with signs $(-,-,+,+)$, and define
\[
 \mathcal L_e^{(s_0,s_1)}=\mathcal R\log(\omega_c\omega_i).
\]
The endpoint-first estimate is
\begin{equation}
\boxed{\begin{gathered}
 |\mathcal L_e^{(k,j)}-\mathcal L_e^{(1,1)}|
 \le R_e\sum_N w_{j,N-v-q'_j}=O_{\rm actual}(q_k),\\
 \limsup_{k\to\infty}\frac{|\mathcal L_e^{(k,j)}-\mathcal L_e^{(1,1)}|}{q_k}
 \le32\log(2+\sqrt3).
\end{gathered}}
\label{eq:long}
\end{equation}
% BEGIN EOR RECEIVER COMMON
\paragraph{Sharper common-loss transport (EOR28--29).}
Replace the displayed width by the intersection width in
\eqref{eor:eq:commonfinite}. On the same paired original families this gives
\[
 \limsup\frac{|\mathcal L_e^{(k,k+4)}-\mathcal L_e^{(1,1)}|}{q_k}
 \le16\log(v_*/u_*)<29.311139553141.
\]
This does not evaluate the absolute order-one common loss. Both target gauges,
$Q_b$, $Q_o$ and the original source denominator remain in the receiving identity.
% END EOR RECEIVER COMMON

For the last constant, the two low widths are $o(j)$ and each high width divided by $j$ tends to $2a_*=\log(2+\sqrt3)$. Now $R_e/k\le16+o(1)$ and $j/k\to1$. This proof compares only four endpoint forms; adding a per-row source-order error over $q_j$ rows would lose the result.

The original IRR/CIT rows are unchanged. Explicitly
\[
 f_{a,n}=\frac{\nu_{a,n+v}-\sum_{d=v+1}^{n+v}\binom{n+v}{d}\mu_df_{a,n+v-d}}
                  {\binom{n+v}{v}\mu_v},
 \qquad \nu_{a,n}=\sum_{u,w}(Z_{\exp,j})_{uw,a}\beta_{j,u,w}^{n},
\]
\[
 b_{n+1,m}=i^{-1}(b_{n,m-1}-c'_jb_{n,m})
              -n(n-1+s_1/2)b_{n-1,m},
\]
\[
 t_n=\left(\frac{\sum_m b_{n,m}f_{a,m}}{\sqrt{h_{n,s_1}}}\right)_a,
 \qquad v_n=\left(\frac{\sum_m b_{n,m}(\beta')^m}{\sqrt{h_{n,s_1}}}\right)_{\beta'}.
\]
Each source family uses its own rows with the same original coefficients. For every fixed flag $Y$, the raw frame is $[\operatorname{diag}(V_D,4),\operatorname{diag}(T_D,4)Y]$. Its positive update retains the complete old Gram and both lower-evaluation and further-quotient Schur terms. Hence the original long weights $1,2,\ldots,2,1$ and the signed difference $\gamma_{H,n}-\gamma_{A,n}$ telescope to exactly the four endpoint forms used above. Neither positive term is discarded, and \eqref{eq:long} does not assign a sign or value to their absolute long sum.

\section{Residual boundary and original arithmetic receivers}
The RQB/CIT identity is retained separately in each source pair:
\begin{equation}
 \mathcal V_b+\mathcal V_{\rm inv}+\mathcal V_o
       =4\widehat F_j^{(s_1)},\qquad
 \mathcal V_b,\mathcal V_{\rm inv},\mathcal V_o\ge0,
 \qquad \mathcal L_e=P_e^{\rm rest}-\mathcal V_b.
 \label{eq:RQB}
\end{equation}
Each term is the same original fixed flag return. For a rank-$r$ target flag form $Q_X$, \eqref{eq:dualord} gives the finite endpoint bound
\[
 -r\log\kappa_j-r u_{j,a_N}
 \le \log\frac{\det Q_X^{(j)}(N)}{\det Q_X^{(1)}(N)}
 \le-r\log\kappa_j-r l_{j,a_N}.
\]
For its dual-sign long return the resulting lower bound is $r(\sum_{\rm low}l-\sum_{\rm high}u)$ and the upper bound is $r(\sum_{\rm low}u-\sum_{\rm high}l)$. The source scalar cancels, and the actual ranks of $Q_b,H_{\rm inv},Q_o$ are retained.

In particular the full source denominator of $Q_b$ is not replaced by a Euclidean norm. It is the original minimum on $V_b/(\im J_0+\ker\Phi_b)$, and the target is the minimum on $W/W_i$. At a displaced source cutoff $N_s\le2q_{k+4}$, its original polynomial quotient degree is at most $2q_{k+4}-v-q'_k=k^2+34k+1-v$. For $k\ge37$ this is at most $2q_k$, so the same full-form theorem also bounds that source minimum. The exact fixed congruence between the CIT and RQB value frames cancels only in matched ratios or the four-return. Thus \eqref{eq:RQB} and CIT's $P_e^{\rm rest}$ retain the source denominator, $A_{\rm inv}$, $Q_o$, and both target gauge losses. The present theorem controls their source-order change, not their separate leading values.

For a proved interval $[L_{X,1},U_{X,1}]$ for an order-one native allocation, add the explicit two endpoints of \eqref{eq:allocationinterval} to obtain an order-$k$ interval; intersect it with all existing intervals. Couple the pair by the exact equality $S_R=\mathcal T-S_K$. No independence of the two error terms is assumed. The original arithmetic receiving identities remain
\begin{equation}
\begin{split}
 F_K^{(s)}&=S_{K,s}-\mathcal E_s^{\rm flag}+F_L^{(s)}-F_{W_0}^{(s)},\\
 \Phi_{k,s}&=F_{W_0}^{(s)}+S_{R,s}+\delta_{0,k,s},\\
 \Psi_k&=F_K^{(k)}+\mathcal H_k+\delta_k^\sigma,
 \qquad J_k^{\rm action}=\Psi_k+C_k^{\rm budget}.
\end{split}\label{eq:return}
\end{equation}
The unequal flag endpoints, the full CEP/$\alpha_A$ transfer, both low metrics, the arithmetic source comparison, and the correlated mixed pairing remain. In particular the matching residual in $C_k^{\rm budget}$ still cancels the same term in $\Psi_k$. A sharper source-order comparison of a kernel volume alone does not lower the exactly combined untilted action allowance.

\section{Exact constants, verification, and remaining calculation}
The constants in the two new limiting allowances satisfy
\begin{equation}
\begin{gathered}
 10.535663175398<8\log(2+\sqrt3)<10.535663175399,\\
 42.142652701594<32\log(2+\sqrt3)<42.142652701595,\\
 15.090354888959<4+16\log2<15.090354888960.
\end{gathered}\label{eq:numerics}
\end{equation}
They have an elementary rational certificate. Enclose $\sqrt3$ between
\[\frac{1732050807568877293527446341505}{10^{30}}\] and that rational plus $10^{-30}$, checking their squares. Use
\[
 \log(2+\sqrt3)=\frac2{\sqrt3}\sum_{n\ge0}\frac{3^{-n}}{2n+1},
\]
whose tail after $m$ terms is at most $3^{-m}/[(2m+1)(1-1/3)]$. Fifty terms suffice. For $\log2$ use $2\sum_{n\ge0}(1/3)^{2n+1}/(2n+1)$ and its geometric tail. All endpoints in \eqref{eq:numerics} follow by rational arithmetic. These are constants, not certified actual period inputs.

The new results are the full native order transport, the explicit $O(q_k)$ allocation and common-long-loss order allowances, the evaluated low-endpoint $q_k$ coefficient, and the $o(q_k)$ paired primitive-to-connecting-image and compatibility differences. They use the complete original fibre maps. The older source-band and total first-variation theorems are not claimed as new results.

The original order-one common long loss, the signed combination $P_e^{\rm rest}$ with its remaining quotients, and the individual $kq_k$-scale coefficients of $S_K,S_R$ remain uncomputed. The present result removes an order-transport obstruction: a subsequent order-one leading interval transfers to the original paired orders with a now specified lower-order error. It does not itself furnish an unconditional narrower numerical leading interval for an individual allocation, an RH exclusion, or a certified offcritical zero. No independent mathematical review or Lean certification is asserted. Execution and PDF inspection records, when delivered, have their separate declared scope.


% BEGIN COMPLETE EOR PROVIDER
\clearpage
\begingroup
\setcounter{equation}{0}
\renewcommand{\theequation}{EOR\arabic{equation}}
\section{Original quantities, inherited inputs and exact comparison domains}
Fix the actual hypothetical simple quartet, the full arithmetic unit, one original
five-orbit period and its logarithm branch. Throughout
\begin{equation}
\begin{gathered}
 k=4l+1\ge29,\quad q=(k+1)^2,\quad q'=(k-7)^2,\quad
 \Delta=16k-48,\quad c'=k/2-4,\\
 \beta_{a,b}^{(n)}=n/2+(2a-n)\delta+i(2b-n)\gamma,
 \quad0<\delta<1/2,\quad\gamma>2,\\
 \chi_n(S)=\prod_{a,b=0}^n(S-\beta_{a,b}^{(n)}),\quad
 \psi=\chi_{k-8},\quad E_A(z)=\sum_{a,b=0}^8a_{ab}e^{\beta_{a,b}^{(8)}z},\\
 v=\operatorname{ord}_0E_A,\quad\mu_v=E_A^{(v)}(0)\ne0,
 \quad\alpha_A=v!/\mu_v,\quad\mathfrak m=\Delta-v.
\end{gathered}\label{eor:eq:data}
\end{equation}
The $a_{ab}$ are the actual fixed-period coefficients, not free numerical inputs.
Their nonreduced conductor is unchanged; the exact $v$ occurs in all cutoffs.
The fixed bound $v\le80$ suffices. No higher-multiplicity metric statement is claimed.

The two source orders are exactly $s=1,k$ with
\begin{equation}
\begin{gathered}
 dm_s(y)=\frac{M_s2^{s/2}}{4\pi\Gamma(s/2)}
 |\Gamma(s/4+iy/2)|^2dy,\quad M_s=(2\pi)^{s/2},\\
 h_{n,s}=M_sn!(s/2)_n,\qquad \vartheta=\pi/2.
\end{gathered}\label{eor:eq:sources}
\end{equation}
Retain the complete original relation matrix $\mathsf U_N$, whose columns are
$\alpha_A\mathcal T_A(\chi_kS^d)/\psi$ for $0\le d\le N-q$,
where $\mathcal T_Ap(S')=\sum a_{ab}p(S'+\beta_{a,b}^{(8)})$.
The increasing low-monomial inclusion is $E$; $Z_k$ and $P_Y$ are the original
kernel frame and residual quotient map. Define the same attained forms
\begin{equation}
\begin{split}
 z^*H_{Y,N}^{(s)}z&=\min_h\int_\R
 |(Ez+\mathsf U_Nh)(c'+iy)|^2|\psi(c'+iy)|^2dm_s(y),\\
 H_{K,N}^{(s)}&=Z_k^*H_{Y,N}^{(s)}Z_k,\qquad
 H_{R,N}^{(s)}=[P_Y(H_{Y,N}^{(s)})^{-1}P_Y^*]^{-1}.
\end{split}\label{eor:eq:forms}
\end{equation}
Both minima are over the full original affine fibres. At the four endpoints
\begin{equation}
\begin{gathered}
 (N_0,N_1,N_2,N_3)=(q-1,q,2q-1,2q),\quad
 (\sigma_0,\sigma_1,\sigma_2,\sigma_3)=(+,+,-,-),\\
 a_N=N-v-q',\quad t_0=\dim(K\cap\C[S]_{<v}),\\
 t_K=8k-16-t_0,\quad t_R=8k-32-v+t_0,\\
 S_X(k;s)=\sum_{i=0}^3\sigma_i\log\det H_{X,N_i}^{(s)},
 \quad\Delta S_X=S_X(k;k)-S_X(k;1),\quad X=K,R.
\end{gathered}\label{eor:eq:returns}
\end{equation}
Zero-dimensional determinants are one. The notation $\Delta S_X$ is a source-order
difference, not the root-strip dimension $\Delta$.

We use GPA's full-form theorem, not a determinant consequence. Its finite domain is
\begin{equation}
 R_k=k\sqrt{\delta^2+\gamma^2}\le2^{-33}q,\qquad k\ge29.
 \label{eor:eq:GPA-domain}
\end{equation}
The new high-degree localization has additional explicit guards in
Section~\ref{eor:sec:localization}. All hold eventually for each fixed stipulated
packet. The old finite QOT bounds remain when a new guard is not satisfied.
Arithmetic maps retain their original period domain $|u|\ge R_*$ or the supplied
exact period tests. No moving-period uniformity is inferred.

\subsection{The complete inherited finite logarithmic constants}
For later use reproduce the constants of GPA/QOT. Let $\epsilon=2^{-32}$,
$h_q=1/q$, $\ell_s=(s-1)/4$, $Q_s=q'+\ell_s$, and
\begin{equation}
\begin{gathered}
 \beta_s=\frac{M_s}{\sqrt2\Gamma(s/2)},\quad\beta_1=1,
 \quad C_L=\pi/\Gamma(3/4)^2,\\
 K_q=\frac{\sqrt{2\pi}}{C_L}\frac{(2q+1)^2}{q}e^{-2q},\qquad
 \mathcal E_k=\frac{qR_k^2}{\epsilon^2q^2-R_k^2}+\log(1+K_q),\\
 r_{\ell_s}=\frac{\sum_{a=0}^{\ell_s-1}(2a+1/2)^2}{\epsilon^2q^2}
                  +\log(1+K_q)\ (\ell_s>0),\qquad r_0=0,\\
 L_h=C_L(\sin h_q)^{3/2},\quad
 U_h=\frac{\Gamma(1/4)^2}{2\pi}(\sin h_q)^{-1/2},\\
 K_{a,Q}=2Q+1+2a+2\sqrt{a(a+2Q)},\\
 b_{k,a,s}^{-}=-\mathcal E_k+\log L_h-K_{a,Q_s}\log(1+h_q/\vartheta),\\
 b_{k,a,s}^{+}=\mathcal E_k+r_{\ell_s}+\log U_h
                       +K_{a,Q_s}\log(1/(1-h_q/\vartheta)).
\end{gathered}\label{eor:eq:GPA-constants}
\end{equation}
Write
$\norm p_{Q}^2=\int_\R|p(y)|^2|y|^{2Q}e^{-\vartheta|y|}dy$.
For every original polynomial of degree at most $a$ in the proved degree domain,
GPA gives
\begin{equation}
 \beta_s e^{b_{k,a,s}^{-}}\norm{p(c'+i\,\cdot)}_{Q_s}^2
 \le\norm{\psi p}_{s,c'}^2
 \le\beta_s e^{b_{k,a,s}^{+}}\norm{p(c'+i\,\cdot)}_{Q_s}^2.
 \label{eor:eq:GPA}
\end{equation}
These are bounds on the entire polynomial norm. The original inside-region error,
root pairs and source mass are in \eqref{eor:eq:GPA-constants}; none is reinterpreted as
a small inverse. At the degrees used here $b^\pm=O_{\rm actual}(\log q)$.
This remains the inherited analytic input, not a new proof attribution.

\section{The original equilibrium and an explicit exterior gap}
Use the proved EIQ equilibrium for
$V_{\alpha,\beta}(x)=\beta\sqrt x-\alpha\log x$ on $(0,\infty)$.
For clarity denote its squared support endpoints by $A,B$ and its density by $\varrho$.
Thus
\begin{equation}
\begin{gathered}
 A=u^2,\quad B=w^2,\quad
 uK(\kappa)=\alpha\pi/\beta,\quad
 wE(\kappa)=(\alpha+2)\pi/\beta,\quad\kappa^2=1-u^2/w^2,\\
 \mathcal P(z)=\int_A^B\log|z-t|\varrho(t)dt,\qquad
 D(x)=V(x)-2\mathcal P(x)-\lambda\ge0,\quad D|_{[A,B]}=0,\\
 \mathfrak h(x)=\frac{\beta}{2\pi x}\int_0^\infty
   \frac{\sqrt{s}\,ds}{(x+s)\sqrt{(A+s)(B+s)}}>0,\\
 D'(x)=\sqrt{(x-A)(x-B)}\,\mathfrak h(x),\quad x\notin[A,B].
\end{gathered}\label{eor:eq:EIQ}
\end{equation}
The last square root is negative on $0<x<A$ and positive on $x>B$.
These are EIQ10--19's actual density and resolvent, including the strict left
exterior inequality. EIQ3--6 and 29 prove smooth positive endpoint dependence
on the positive parameter quadrant.

For each such parameter pair define explicit positive constants
\begin{equation}
\begin{gathered}
 g=B-A,\qquad
 h_* =\frac{\beta}{2\pi(B+1)(B+2)(B+3)},\qquad
 c_* =\frac23h_*\sqrt g,\\
 \kappa_* =\min\left\{A h_*\sqrt g,
               \frac{\beta\sqrt B}{6\pi(3B+2)}\right\},\qquad
 \eta_* =\min\{A/2,1\},\\
 L=\max_{x\in[A,B]}|\beta/(2\sqrt x)-\alpha/x|,
 \qquad C_*^2=32e^{3L/2}.
\end{gathered}\label{eor:eq:gapconstants}
\end{equation}
The stars here label local gap constants, not any canonical allowance from the
programme. The maximum $L$ is an elementary finite maximum: evaluate the absolute
expression at $A,B$ and at $(4\alpha/\beta)^2$ if that point belongs to $[A,B]$.

\begin{lemma}[Both exterior tails, with a global upper derivative bound]
For $0<\eta\le\eta_*$,
\begin{equation}
\begin{aligned}
 D(x)&\ge c_*\eta^{3/2}
        &&(0<x\le A-\eta\ \hbox{or}\ x\ge B+\eta),\\
 xD'(x)&\ge\kappa_*\sqrt\eta&& (x\ge B+\eta).
\end{aligned}\label{eor:eq:gap}
\end{equation}
\end{lemma}
\begin{proof}
In the integral in \eqref{eor:eq:EIQ}, restrict $s$ to $[1,2]$. For
$x\in[A/2,B+1]$, $x\le B+1$, $x+s\le B+3$ and
$\sqrt{(A+s)(B+s)}\le B+2$, whereas $\sqrt s\ge1$.
Thus $\mathfrak h(x)\ge h_*$. On either endpoint segment of length $\eta$,
$|D'(x)|\ge h_*\sqrt g\sqrt{\operatorname{dist}(x,[A,B])}$.
Integrating proves the first assertion at the segment endpoints. Monotonicity
from \eqref{eor:eq:EIQ} proves it on both entire exterior tails.

For $B+\eta\le x\le B+1$ the same bound gives
$xD'(x)\ge A h_*\sqrt g\sqrt\eta$.
For $x\ge B+1$ restrict the integral instead to $s\in[x,2x]$ to obtain
\[
 \mathfrak h(x)\ge\frac{\beta}{6\pi\sqrt x(B+2x)},\qquad
 xD'(x)\ge\frac{\beta\sqrt x(x-B)}{6\pi(B+2x)}
 \ge\frac{\beta\sqrt B}{6\pi(3B+2)}.
\]
The last ratio is minimized at $x=B+1$ after the factor $\sqrt x$ is bounded
below by $\sqrt B$. Since $\sqrt\eta\le1$, this proves the second assertion
with the displayed common constant. All inequalities concern the same EIQ
potential and its complete positive half-line.
\end{proof}

\section{Uniform polynomial localization and a growing moment shift}
\label{eor:sec:localization}
Let $m$ be a positive integer and $P$ any complex polynomial of degree at most $m$.
Set $\norm P_{m,V}^2=\int_0^\infty|P(x)|^2e^{-mV(x)}dx$.
The following complete evaluation bound is the argument of OOQ7, restated to
specify its constants for the parameter pair used here:
\begin{equation}
 |P(x)|^2e^{-mV(x)}\le C_*^2m^4e^{-mD(x)}\norm P_{m,V}^2,
 \quad x\notin[A,B],\quad m\ge\max\{1,2/g\}.
 \label{eor:eq:eval}
\end{equation}
For completeness, let $M=\max_{[A,B]}|P|e^{-mV/2}$, attained at $x_0$.
One side of $x_0$ has length $1/m$ inside the support. There
$|P(x)|\le e^{L/2}|P(x_0)|$.
A Chebyshev expansion on an interval $I$ proves
$\norm{P'}_{\infty,I}\le4m^3|I|^{-1}\norm P_{\infty,I}$:
the nonconstant coefficients have modulus at most $2\norm P_\infty$ and
$|T_j'|\le j^2$, so the sum and affine derivative factor give the bound.
It follows that $|P|\ge|P(x_0)|/2$ on a segment of length
$(8e^{L/2}m^4)^{-1}$. The weight there is at least $e^{-L}e^{-mV(x_0)}$.
Hence $\norm P_{m,V}^2\ge M^2/(32e^{3L/2}m^4)$.
Finally $\log|P(z)|-m\mathcal P(z)$ is subharmonic off the support,
where $\mathcal P$ is the logarithmic potential in \eqref{eor:eq:EIQ}.
The equilibrium equality bounds its boundary values by $\log M+m\lambda/2$.
It extends subharmonically at infinity, since $\deg P\le m$.
The maximum principle proves \eqref{eor:eq:eval}. Zeros give only negative singularities.
This uses the logarithmic potential, not a partition-function error estimate.

Define
\begin{equation}
\begin{gathered}
 \eta_m=\left(
 \frac{\log(C_*^2(B+2))+10\log(m+1)}{c_*m}
 \right)^{2/3},\qquad \varepsilon_m=(m+1)^{-6},\\
 m\ge\max\{1,2/g\},\qquad
 \eta_m\le\eta_*,\qquad
 m\kappa_*\sqrt{\eta_m}\ge\ell+2 .
\end{gathered}\label{eor:eq:guards}
\end{equation}
The last line gives all additional finite guards; $\ell\ge0$ may be real.

\begin{theorem}[Uniform norm comparison from the actual support]
Under \eqref{eor:eq:guards}, every degree-at-most-$m$ complex polynomial satisfies
\begin{equation}
 \boxed{\quad
 (A-\eta_m)^\ell(1-\varepsilon_m)\norm P_{m,V}^2
 \le\int_0^\infty x^\ell|P(x)|^2e^{-mV(x)}dx
 \le(B+\eta_m)^\ell(1+\varepsilon_m)\norm P_{m,V}^2.
 \quad}\label{eor:eq:localization}
\end{equation}
\end{theorem}
\begin{proof}
For any $0<\eta\le\eta_*$, \eqref{eor:eq:eval} and the left gap give
\[
 \frac{\int_0^{A-\eta}|P|^2e^{-mV}dx}{\norm P_{m,V}^2}
 \le C_*^2m^4 A e^{-mc_*\eta^{3/2}}.
\]
Consequently the numerator in \eqref{eor:eq:localization} is at least
$(A-\eta)^\ell$ times one minus this mass fraction.
For the upper tail integrate the derivative inequality in \eqref{eor:eq:gap}:
\[
 D(x)\ge c_*\eta^{3/2}+\kappa_*\sqrt\eta\log\frac{x}{B+\eta},
 \qquad x\ge B+\eta.
\]
If $m\kappa_*\sqrt\eta>\ell+1$, direct integration gives
\[
 \int_{B+\eta}^\infty x^\ell e^{-mD(x)}dx
 \le\frac{(B+\eta)^{\ell+1}e^{-mc_*\eta^{3/2}}}
                {m\kappa_*\sqrt\eta-\ell-1}.
\]
The portion up to $B+\eta$ is bounded by $(B+\eta)^\ell\norm P_{m,V}^2$;
use \eqref{eor:eq:eval} on the rest. At the chosen $\eta_m$,
\[
 e^{-mc_*\eta_m^{3/2}}
   =[C_*^2(B+2)(m+1)^{10}]^{-1}.
\]
The lower omitted mass and the upper added fraction, after division by
$(B+\eta_m)^\ell$, are both at most $(m+1)^{-6}$:
$A<B+2$, $B+\eta_m<B+2$, $m^4\le(m+1)^4$, and the upper denominator is at
least one. This proves both full integral inequalities, without assuming that
$P$ avoids either tail.
\end{proof}

On any compact positive parameter rectangle the constants in
\eqref{eor:eq:gapconstants} have uniform positive lower bounds and finite upper bounds.
This follows from the proved smooth positive endpoints; it is not a condition
on an unspecified polynomial. Thus, uniformly for $\ell=O(\sqrt m)$,
\begin{equation}
 \eta_m=O((\log(m+1)/m)^{2/3}),\quad
 \ell\eta_m=o(1),\quad \varepsilon_m=O(m^{-6}).
 \label{eor:eq:locasym}
\end{equation}
All guards hold eventually. In particular the logarithmic endpoints in
\eqref{eor:eq:localization} are $\ell\log A+o(1)$ and $\ell\log B+o(1)$.
These are uniform enclosures; equality with either endpoint has not been asserted.

\section{Both square-map sheets and the finite improved native interval}
\label{eor:sec:application}
Apply the localization to the high polynomial degrees in \eqref{eor:eq:returns}.
For a general polynomial $p_c(y)=p(c'+iy)$ of degree at most $a$, write exactly
\begin{equation}
 p_c(y)=P_0(y^2/q^2)+(y/q)P_1(y^2/q^2),\quad
 m_0=\lfloor a/2\rfloor,\quad m_1=\lfloor(a-1)/2\rfloor .
 \label{eor:eq:parity}
\end{equation}
This defines both $P_\epsilon$ from the same original coefficients, with no
reality assumption. The odd cross terms integrate to zero in the even reference
weight, so
\begin{equation}
 \norm{p_c}_{q'}^2=q^{2q'+1}\sum_{\epsilon=0}^1
 \int_0^\infty |P_\epsilon(x)|^2
 x^{q'+\epsilon-1/2}e^{-\vartheta q\sqrt x}dx.
 \label{eor:eq:square}
\end{equation}
The factor $q^{2q'+1}$ includes both half-lines and their Jacobians.
After shifting $q'$ to $q'+\ell$, the right side has the extra factor
$q^{2\ell}x^\ell$. Both parity blocks remain.

For each $\epsilon=0,1$ set
\begin{equation}
 \alpha_\epsilon=\frac{q'+\epsilon-1/2}{m_\epsilon},\qquad
 \beta_\epsilon=\frac{\vartheta q}{m_\epsilon},\qquad
 [A_\epsilon,B_\epsilon]=\operatorname{supp}\mu_{\alpha_\epsilon,\beta_\epsilon}.
 \label{eor:eq:parameters}
\end{equation}
Use the localization only when $m_\epsilon\ge1$, both parameter pairs belong to
$[3/2,5/2]\times[\pi/2,3\pi/2]$, and \eqref{eor:eq:guards} holds in both blocks.
These tests are explicit in the actual integers and the unique EIQ endpoints.
For the two original high cutoffs they hold eventually. Below that domain keep
QOT's finite Laguerre bounds unchanged.

Let $\eta_\epsilon,\varepsilon_\epsilon$ be \eqref{eor:eq:guards}'s values. Define
\begin{equation}
\begin{split}
 F_{k,a}^{-}&=\min_{\epsilon=0,1}
  \{\ell\log(A_\epsilon-\eta_\epsilon)+\log(1-\varepsilon_\epsilon)\},\\
 F_{k,a}^{+}&=\max_{\epsilon=0,1}
  \{\ell\log(B_\epsilon+\eta_\epsilon)+\log(1+\varepsilon_\epsilon)\},\\
 L_{k,a}^{E}&=b_{k,a,k}^{-}-b_{k,a,1}^{+}
       +2\ell\log\frac{q\vartheta}{2q'}+F_{k,a}^{-},\\
 U_{k,a}^{E}&=b_{k,a,k}^{+}-b_{k,a,1}^{-}
       +2\ell\log\frac{q\vartheta}{2q'}+F_{k,a}^{+},\\
 \kappa_k&=\beta_k(4q'/\pi)^{2\ell}
          =2^{10\ell}\frac{(2\ell)!}{(4\ell)!}(k-7)^{4\ell}.
\end{split}\label{eor:eq:newconstants}
\end{equation}
The comparison centre is exactly the earlier rational $\kappa_k$; changing the
bounding endpoints does not change its meaning or absorb a mass.

\begin{theorem}[Full high-endpoint source-order comparison]
For every degree-at-most-$a$ original polynomial in the stated joint domain,
\begin{equation}
 \kappa_k e^{L_{k,a}^{E}}\norm{\psi p}_{1,c'}^2
 \le \norm{\psi p}_{k,c'}^2
 \le\kappa_k e^{U_{k,a}^{E}}\norm{\psi p}_{1,c'}^2.
 \label{eor:eq:newform}
\end{equation}
It holds after every original fixed restriction or attained quotient, in
particular all three forms in \eqref{eor:eq:forms} at the actual cutoff $a=a_N$.
On their original complex-linear dual metric it is reversed:
\begin{equation}
 \kappa_k^{-1}e^{-U_{k,a}^{E}}D^{(1)}
 \preceq D^{(k)}\preceq
 \kappa_k^{-1}e^{-L_{k,a}^{E}}D^{(1)}.
 \label{eor:eq:newdual}
\end{equation}
\end{theorem}
\begin{proof}
Apply \eqref{eor:eq:localization} separately to the exact two summands in
\eqref{eor:eq:square}. Taking the minimum lower constant and maximum upper constant
yields
$e^{2\ell\log q+F^-}\norm{p_c}_{q'}^2
\le\norm{p_c}_{q'+\ell}^2
\le e^{2\ell\log q+F^+}\norm{p_c}_{q'}^2$.
Combine the lower inequality with the lower order-$k$ and upper order-one sides
of \eqref{eor:eq:GPA}, and combine the opposite three inequalities for the upper side.
This is exactly \eqref{eor:eq:newconstants}--\eqref{eor:eq:newform}.
For each affine fibre every candidate lift satisfies the two inequalities;
taking both infima preserves them. For a restriction substitute the full fixed
columns. These operations keep all relation, kernel and residual mixed blocks.
Positive-form inversion reverses the order, and the source-independent congruence
$D=A_{\ell,k}^*\overline{H_Y^{-1}}A_{\ell,k}$ gives \eqref{eor:eq:newdual}.
It too passes to the same fixed target flag subquotients by the identical
minimum argument. No norm is imposed on an observed coordinate space.
\end{proof}

For a guaranteed intersection with the earlier finite bounds reproduce QOT's
endpoints. For $\alpha_k=2q'$ put
\[
 d^-_{k,a}=4\sum_{r=0}^{\ell-1}\arsinh\sqrt{(a+r+1/2)/\alpha_k},\qquad
 d^+_{k,a}=4\sum_{r=1}^{\ell}\arsinh\sqrt{(a+r+1/2)/\alpha_k},
\]
\[
 L_{k,a}^{Q}=b^-_{k,a,k}-b^+_{k,a,1}-d^-_{k,a},\qquad
 U_{k,a}^{Q}=b^+_{k,a,k}-b^-_{k,a,1}+d^+_{k,a}.
\]
At an endpoint where the new guards hold, use
\begin{equation}
 \underline b_{k,a}=\max\{L_{k,a}^Q,L_{k,a}^E\},\qquad
 \overline b_{k,a}=\min\{U_{k,a}^Q,U_{k,a}^E\}.
 \label{eor:eq:intersection}
\end{equation}
Else use the QOT endpoints alone. The new finite enclosure never discards a
stronger earlier one. QOT's low bounds remain part of this definition.

\section{Evaluated asymmetric constants at scale $q_k$}
Let $u_*,v_*$ denote the original unsquared EIQ endpoints at $(2,\pi)$:
\begin{equation}
 u_*K(\kappa)=2,\qquad v_*E(\kappa)=4,\qquad
 \kappa^2=1-u_*^2/v_*^2.
 \label{eor:eq:limitendpoints}
\end{equation}
They are not period variables, and no independent Gamma order is introduced.
At either high cutoff $a_N=q+\Delta-v+O(1)$, so
\[
 m_\epsilon=q/2+O(k),\qquad
 (\alpha_\epsilon,\beta_\epsilon)=(2,\pi)+O(k^{-1}).
\]
Smooth endpoint dependence, \eqref{eor:eq:locasym}, and
$b^\pm=O_{\rm actual}(\log q)$ now give
\begin{equation}
\begin{aligned}
 L_{k,a_N}^{E}&=k a_-+O_{\rm actual}(\log q),
      &a_-&=\tfrac12\log(\pi u_*/4),\\
 U_{k,a_N}^{E}&=k a_++O_{\rm actual}(\log q),
      &a_+&=\tfrac12\log(\pi v_*/4).
\end{aligned}\label{eor:eq:highasym}
\end{equation}
For example $F^-=(k/4)\log(u_*^2)+O(1)$ because
$\ell\eta_{m_\epsilon}=o(1)$ and the endpoint displacement costs
$O(\ell/k)=O(1)$. The remaining centre is
$2\ell\log(\pi q/(4q'))=(k/2)\log(\pi/4)+O(1)$.
The upper proof is identical. These are asymptotics of proved bounding constants,
not a claim about the unknown eigenvalues of the actual high compressed forms.

Taking rank-$t_X$ determinants and the four original signs in
\eqref{eor:eq:newform} yields the finite interval
\begin{equation}
\begin{split}
 t_X(\underline b_{k,a_{N_0}}+\underline b_{k,a_{N_1}}
     -\overline b_{k,a_{N_2}}-\overline b_{k,a_{N_3}})
 &\le\Delta S_X\\
 &\le
 t_X(\overline b_{k,a_{N_0}}+\overline b_{k,a_{N_1}}
     -\underline b_{k,a_{N_2}}-\underline b_{k,a_{N_3}}).
\end{split}\label{eor:eq:allocationfinite}
\end{equation}
The four identical rank factors multiplying $\log\kappa_k$ cancel exactly here,
not before the endpoint forms are defined.
At a low endpoint $a_N=O(k)$ the retained QOT endpoints are
$O_{\rm actual}(\sqrt k+\log q)$; this follows directly from
$\arsinh x\le x$ in its sums. Since $t_X=8k+O(1)$,
\begin{theorem}[Narrower source-order allocation interval]
For either original allocation $X=K,R$ at a fixed stipulated packet and period,
\begin{equation}
 \boxed{\quad
 -C_+\le\liminf_{k\equiv1\ (4)}\frac{\Delta S_X}{q_k}
 \le\limsup_{k\equiv1\ (4)}\frac{\Delta S_X}{q_k}\le C_-,
 \quad}\label{eor:eq:alloclimit}
\end{equation}
where
\begin{equation}
 C_+=8\log(\pi v_*/4),\qquad
 C_-=8\log(4/(\pi u_*)).
 \label{eor:eq:Cpm}
\end{equation}
More explicitly the finite interval \eqref{eor:eq:allocationfinite} implies
$-C_+q_k-O_{\rm actual}(k^{3/2}+k\log q_k)\le\Delta S_X
\le C_-q_k+O_{\rm actual}(k^{3/2}+k\log q_k)$.
No unique normalized limit is asserted.
\end{theorem}
\begin{proof}
Use \eqref{eor:eq:highasym} at the two high endpoints and the preceding low bound.
The lower centre is $-2t_Xka_+$, the upper centre is $-2t_Xka_-$.
Replacing $t_Xk$ by $8q_k$ costs only $O(k)$ because $t_0,v$ are bounded.
Substitution of $a_\pm$ gives the stated constants and finite order of the
remaining allowance. The explicit finite endpoints, rather than an unspecified
big-$O$ constant, are available in \eqref{eor:eq:allocationfinite}.
\end{proof}

\section{The original common long loss and its fixed-frame receiver}
Use $j=k+4$ and the original paired source orders $(1,1)$ and $(k,j)$.
In target coordinates the source-independent metric is
$\mathcal D_{j,N}=\operatorname{diag}(D_{j,N-v},4)$.
The complete original flags are
\begin{equation}
\begin{gathered}
 I=\im(B_{\rm cof}J_0),\quad B_0=\im B_{\rm cof},\quad
 S=\im[B_{\rm cof},F_e],\quad J=\ker\mathbb B,\\
 K_i=I\cap J,\quad W_i=I+J,\quad W=B_0+J,\quad
 K_c=S\cap W,\quad T=S+J,\\
 \omega_i=\frac{\det\Gr(W_i/J)}{\det\Gr(I/K_i)},\quad
 \omega_c=\frac{\det\Gr(T/W)}{\det\Gr(S/K_c)},\quad
 R_e=\dim(I/K_i)+\dim(S/K_c)\le16k-8.
\end{gathered}\label{eor:eq:flags}
\end{equation}
All Grams are their attained native quotient forms in matched value frames.
The fixedness of $F_e$ under the source presentation follows from the same
Taylor cochain equation and uniqueness modulo all lower evaluation columns,
as in QOT's full primitive map. No corresponding denominator or multiplier
kernel is changed by the present scalar comparison.

Let $\mathcal R$ have signs $(-,-,+,+)$ at
$N=q_j-1,q_j,2q_j-1,2q_j$ and write
$\mathcal L_e=\mathcal R\log(\omega_c\omega_i)$.
With $w_{j,a}=\overline b_{j,a}-\underline b_{j,a}$, subtraction of the
matched quotient determinant intervals in \eqref{eor:eq:newdual} proves
\begin{equation}
 \boxed{\quad
 |\mathcal L_e^{(k,j)}-\mathcal L_e^{(1,1)}|
       \le R_e\sum_Nw_{j,N-v-q'_j}.
 \quad}\label{eor:eq:commonfinite}
\end{equation}
The common factor $\kappa_j^{-r}$ cancels in each rank-$r$ quotient contraction
at each endpoint. Thus a small observation inverse is not needed. This estimate
is applied at four full endpoint forms, not summed as an independent error over
$q_j$ rows.
By \eqref{eor:eq:highasym}, each high width equals
$\frac j2\log(v_*/u_*)+O_{\rm actual}(\log q_j)$, and both low widths are
$O_{\rm actual}(\sqrt j+\log q_j)$. Hence
\begin{equation}
 \boxed{\quad
 \limsup\frac{|\mathcal L_e^{(k,k+4)}-\mathcal L_e^{(1,1)}|}{q_k}
 \le C_{\rm com}:=16\log(v_*/u_*)=2(C_++C_-).
 \quad}\label{eor:eq:commonlimit}
\end{equation}
The error about this bound is $O_{\rm actual}(k^{3/2}+k\log q_k)$,
with the complete finite version \eqref{eor:eq:commonfinite}.
This does not replace the earlier one-budget estimate for the absolute
order-one common long loss by an evaluated value.

\subsection{The actual rows, half-lines and all remaining terms}
The unchanged source-coordinate pullback is
$(\mathsf P_{c',a})_{r,n}=\binom nr(c')^{n-r}i^r$.
The full power-exponential reference rows remain
\begin{equation}
\begin{gathered}
 \mathsf A_{b,r}=(-1)^b\binom rb\Gamma(r+2Q_s+1)
 \sqrt{\frac{b!}{\Gamma(b+2Q_s+1)}}\,
 \vartheta^{-r-(2Q_s+1)/2},\\
 \mathsf R_{b,r}=(-1)^b2^{r-b}\binom{2Q_s+r}{r-b}
 \sqrt{\frac{r!\Gamma(b+2Q_s+1)}{b!\Gamma(r+2Q_s+1)}},\\
 \mathsf W=\begin{bmatrix}\mathsf A\mathsf P_{c',a}\\
                           \mathsf R\mathsf A\mathsf P_{c',a}\end{bmatrix}.
\end{gathered}\label{eor:eq:W}
\end{equation}
Zero triangular entries are understood. Formula \eqref{eor:eq:square} is merely a
second exact coefficient presentation of this same complete two-half-line norm.
For every original column $x$, both row blocks of $\mathsf Wx$ remain; only their
integrated even/odd cross terms cancel by symmetry of the reference density.
Original period coefficients are not conjugated by reflection.
The new theorem acts on these full quadratic forms before either Schur minimum.

The native rows are still calculated from the original symbol and Gamma recurrence:
\begin{equation}
\begin{gathered}
 f_{a,n}=\frac{\nu_{a,n+v}-\sum_{d=v+1}^{n+v}\binom{n+v}{d}
               \mu_df_{a,n+v-d}}{\binom{n+v}{v}\mu_v},\quad
 \nu_{a,n}=\sum_{u,w}(Z_{\exp,j})_{uw,a}(\beta_{u,w}^{(j)})^n,\\
 b_{n+1,r}=i^{-1}(b_{n,r-1}-c'_jb_{n,r})
                 -n(n-1+s_1/2)b_{n-1,r},\\
 t_n=\left(\frac{\sum_r b_{n,r}f_{a,r}}{\sqrt{M_{s_1}n!(s_1/2)_n}}\right)_a,
 \quad
 v_n=\left(\frac{\sum_r b_{n,r}(\beta')^r}{\sqrt{M_{s_1}n!(s_1/2)_n}}\right)_{\beta'}.
\end{gathered}\label{eor:eq:rows}
\end{equation}
In a new step the preceding complete lower Gram is $V^*V$ and the corrected row
is $\rho_n=t_n-v_n(V^*V)^{-1}V^*T$.
Its update denominator is $1+v_n(V^*V)^{-1}v_n^*$.
The subsequent boundary or invariant minimum also remains. None of these rows
or denominators is approximated by its rank in \eqref{eor:eq:commonfinite}.

For an order-one interval $[L_{X,1},U_{X,1}]$, add the two finite endpoints in
\eqref{eor:eq:allocationfinite}, intersect with all earlier bounds, and couple by
the exact $S_R=\mathcal T-S_K$. In the common loss add and subtract the explicit
allowance in \eqref{eor:eq:commonfinite}. Retain separately
\begin{equation}
\begin{gathered}
 \mathcal L_e=P_e^{\rm rest}-\mathcal V_b,\qquad
 \mathcal V_b+\mathcal V_{\rm inv}+\mathcal V_o=4\widehat F_j,\\
 F_K^{(s)}=S_{K,s}-\mathcal E_s^{\rm flag}+F_L^{(s)}-F_{W_0}^{(s)},
 \qquad \Phi_{k,s}=F_{W_0}^{(s)}+S_{R,s}+\delta_{0,k,s},\\
 \Psi_k=F_K^{(k)}+\mathcal H_k+\delta_k^\sigma,
 \qquad J_k^{\rm action}=\Psi_k+C_k^{\rm budget}.
\end{gathered}\label{eor:eq:receivers}
\end{equation}
The original source denominator, $A_{\rm inv}$, $Q_b$, $Q_o$, both target gauges,
both low metrics, unequal flag errors and the complete source-transfer correction
are unchanged. The matching residual cancellation in the full action budget
is not altered by narrowing a source-order interval. The previously proved
low primitive $o(q_k)$ paired source/image comparison is retained without
being reattributed as a new theorem.

\section{Certified constants and the precise remaining task}
Let $r=u_*/v_*\in(0,1)$ and define the normalized elliptic integrals
\[
 \bar K=\frac{2}{\pi}K(\sqrt{1-r^2}),\qquad
 \bar E=\frac{2}{\pi}E(\sqrt{1-r^2}).
\]
The endpoint equation is $r\bar K/\bar E=1/2$, and
\begin{equation}
 C_+=8\log(2/\bar E),\qquad C_-=8\log\bar K,
 \qquad C_{\rm com}=-16\log r.
 \label{eor:eq:ellipticconstants}
\end{equation}
Thus no numerical enclosure of $\pi$ is required for these three constants.
The accompanying integer/rational calculation certifies
\begin{equation}
\begin{gathered}
 0.16010167095771835588456745<r<
 0.16010167095771835588456746,\\
 8.881817865596<C_+<8.881817865597,\\
 5.773751910973<C_-<5.773751910974,\\
 29.311139553140<C_{\rm com}<29.311139553141.
\end{gathered}\label{eor:eq:certified}
\end{equation}
Here is the rigorous basis of that certificate. For $m=1-r^2$ and
$c_n=\binom{2n}{n}^2/16^n$,
\[
 \bar K=\sum_{n\ge0}c_nm^n,\qquad
 \bar E=1-\sum_{n\ge1}\frac{c_nm^n}{2n-1}.
\]
These follow by binomial expansion in the defining elliptic integrals and the
integral of $\sin^{2n}\theta$. If $t_n=c_nm^n$, then
$t_{n+1}/t_n=m((2n+1)/(2n+2))^2\le m$.
After degree $N$, the omitted $\bar K$ tail is at most $t_Nm/(1-m)$ and the
omitted $1-\bar E$ tail at most $t_Nm/((1-m)(2N+1))$.
At both rational endpoints in \eqref{eor:eq:certified}, use $N=4096$, enclosing each
recurrence operation by integer floor and ceiling at a fixed denominator
$10^{80}$. The resulting intervals certify opposite signs of
$2r\bar K-\bar E$. EIQ's strict monotonicity then encloses the unique root.
Monotonicity in $m$ encloses both normalized elliptic integrals between these
endpoint evaluations. Finally, for $x>1$ and $z=(x-1)/(x+1)$, use
\[
 \log x=2\sum_{n=0}^{M-1}\frac{z^{2n+1}}{2n+1}+R_M,
 \qquad0<R_M\le\frac{2z^{2M+1}}{(2M+1)(1-z^2)}.
\]
Integer interval arithmetic with $M=400$ proves the displayed decimal brackets.
The full script records its actual rational endpoints and rounding directions.
These are universal-constant certificates, not actual-period evaluations.

The new allocation interval $[-C_+,C_-]$ is strictly contained in the earlier
$[-8\log(2+\sqrt3),8\log(2+\sqrt3)]$ interval.
The new common source-order allowance is strictly below
$32\log(2+\sqrt3)$. These are genuine narrower $q_k$-scale enclosures for the
same differences. They do not assert optimality of the enclosures or existence
of either normalized source-order limit.

What remains is the absolute order-one common long loss and its allocation
between the actual kernel and residual. Neither an individual $kq_k$-scale
coefficient nor a narrower unconditional interval for that coefficient is
assigned here. The current result makes a subsequent order-one estimate more
precisely transferable; it does not substitute for that estimate.
The source comparison is uniform in all polynomial directions on its stated
finite domain; the later arithmetic and geometric maps retain their fixed-period
and simple-quartet hypotheses. No RH conclusion, actual offcritical-zero
certificate, independent referee acceptance, or Lean execution is claimed.


\endgroup
% END COMPLETE EOR PROVIDER

\begin{thebibliography}{9}
\bibitem{gpa}Split-Zero research programme. (2026). \emph{Current kernel and multiplicity proofs} (supplied 498-page edition). Complete GPA1--23, LVM1--19, DNE, SAR, CIT/ITL receiving providers. Original source filename: \texttt{14\_CURRENT\_KERNEL\_AND\_MULTIPLICITY\_PROOFS(6).tex}.
\bibitem{pds}Split-Zero research programme. (2026). \emph{Full primitive connecting proofs}. Complete PDS1--24, CIT1--26 and PDI1--16 providers in the supplied 13-page source; the same-coordinate argument is also recorded in PSR23--24 of the later 40-page supplement.
\bibitem{current}Split-Zero research programme. (2026). \emph{Primitive-band continuation} (549-page edition and 40-page supplement). Repository path \path{workbenches/splitzero-tandem/continuations/20260915-primitive-band-continuation}; consulted at commit \texttt{1efd53337561d8f67cf0ac1d119acdaa222644c3}. Its evaluated totals and later RQB/PSR identities retain their own proof provenance.
\end{thebibliography}
\end{document}
